\chapter{1989年广东卷（文）}

\section{单选题}

\begin{problem}
函数 \(f(x)=x^{\frac{1}{5}}\)（\quad）

\choices
    {是奇函数而不是偶函数}
    {是偶函数而不是奇函数}
    {既是奇函数又是偶函数}
    {既不是奇函数又不是偶函数}
\end{problem}

\begin{answer}
A
\end{answer}

\begin{solution}
函数的定义域是 \(\R\)，且对任意 \(x\in\R\)，有 \(f(-x)=(-x)^{\frac15}=-x^{\frac15}=-f(x)\). 因此 \(f(x)\) 是奇函数. 又例如 \(f(1)=1\neq f(-1)=-1\)，所以它不是偶函数，故选 A.
\end{solution}

\begin{problem}
函数 \(y=\lg x\)（\(x>0\)）的值域是（\quad）

\choices
    {\((0,+\infty)\)}
    {\((-\infty,0)\)}
    {\((-\infty,+\infty)\)}
    {\([0,+\infty)\)}
\end{problem}

\begin{answer}
C
\end{answer}

\begin{solution}
当 \(x>0\) 时，\(\lg x\) 的定义域为 \((0,+\infty)\). 对任意实数 \(y\)，取 \(x=10^y>0\)，则 \(\lg x=y\)，所以函数值可以取遍所有实数，值域为 \((-\infty,+\infty)\)，故选 C.
\end{solution}

\begin{problem}
直线 \(x+\sqrt{3}y-3=0\) 的倾斜角是（\quad）

\choices
    {\(\frac{\pi}{6}\)}
    {\(\frac{\pi}{3}\)}
    {\(\frac{2\pi}{3}\)}
    {\(\frac{5\pi}{6}\)}
\end{problem}

\begin{answer}
D
\end{answer}

\begin{solution}
将直线方程化为斜截式，得 \(y=-\frac{1}{\sqrt{3}}x+\sqrt{3}\)，所以斜率 \(k=-\frac{1}{\sqrt{3}}\). 设倾斜角为 \(\alpha\in[0,\pi)\)，则 \(\tan\alpha=k=-\frac{1}{\sqrt{3}}\)，故 \(\alpha=\frac{5\pi}{6}\)，选 D.
\end{solution}

\begin{problem}
二面角是指（\quad）

\choices
    {两个平面所组成的角}
    {经过同一直线的两个平面所组成的图形}
    {从一条直线出发的两个半平面所组成的图形}
    {两个平面所夹的不大于 \(90^\circ\) 的角}
\end{problem}

\begin{answer}
C
\end{answer}

\begin{solution}
二面角的棱是一条直线，它把两个平面各分成两个半平面；取两个以这条直线为公共边的半平面，所组成的图形叫作二面角. 因此选项 C 正确.
\end{solution}

\begin{problem}
对数方程 \(2\log_{x}8-3\log_{8}x=1\) 的解集是（\quad）

\choices
    {\(\left\{4,\frac{1}{8}\right\}\)}
    {\(\left\{2,\frac{\sqrt{2}}{4}\right\}\)}
    {\(\left\{4,\frac{1}{4}\right\}\)}
    {\(\left\{2,\frac{1}{8}\right\}\)}
\end{problem}

\begin{answer}
A
\end{answer}

\begin{solution}
由对数定义，必须有 \(x>0\) 且 \(x\neq1\). 令 \(t=\log_8x\)，则 \(t\neq0\)，并且 \(\log_x8=\frac1t\). 原方程化为 \(\frac2t-3t=1\)，即 \(3t^2+t-2=0\)，所以 \(t=\frac23\) 或 \(t=-1\). 因而 \(x=8^{\frac23}=4\) 或 \(x=8^{-1}=\frac18\)，解集为 \(\left\{4,\frac18\right\}\)，故选 A.
\end{solution}

\begin{problem}
如果 \(\sin(\pi+A)=-\frac{1}{2}\)，那么 \(\cos\left(\frac{3\pi}{2}-A\right)=\)（\quad）

\choices
    {\(\frac{1}{2}\)}
    {\(-\frac{1}{2}\)}
    {\(\frac{\sqrt{3}}{2}\)}
    {\(-\frac{\sqrt{3}}{2}\)}
\end{problem}

\begin{answer}
B
\end{answer}

\begin{solution}
由 \(\sin(\pi+A)=-\sin A=-\frac12\)，得 \(\sin A=\frac12\). 又
\(\cos\left(\frac{3\pi}{2}-A\right)=\cos\frac{3\pi}{2}\cos A+\sin\frac{3\pi}{2}\sin A=-\sin A=-\frac12\)，故选 B.
\end{solution}

\begin{problem}
平移坐标轴，把原点移动到 \(O'(3,-4)\)，那么点 \((3,-2)\) 的新坐标是（\quad）

\choices
    {\((6,-6)\)}
    {\((-3,4)\)}
    {\((0,-2)\)}
    {\((0,2)\)}
\end{problem}

\begin{answer}
D
\end{answer}

\begin{solution}
设原坐标为 \((x,y)\)，新原点为 \(O'(3,-4)\)，且坐标轴方向不变，则新坐标为 \((x-3,y+4)\). 对点 \((3,-2)\)，新坐标是 \((0,2)\)，故选 D.
\end{solution}

\begin{problem}
如果直线 \(l\) 与直线 \(y=2x-1\) 关于直线 \(y=x\) 对称，那么 \(l\) 的方程是（\quad）

\choices
    {\(y=2x-1\)}
    {\(y=\frac{1}{2}x+1\)}
    {\(y=\frac{1}{2}x-\frac{1}{2}\)}
    {\(y=\frac{1}{2}x+\frac{1}{2}\)}
\end{problem}

\begin{answer}
D
\end{answer}

\begin{solution}
关于直线 \(y=x\) 对称时，点的横、纵坐标互换. 原直线上的点满足 \(y=2x-1\)，交换后满足 \(x=2y-1\)，即 \(y=\frac12x+\frac12\)，故选 D.
\end{solution}

\begin{problem}
设 \(A\) 是直角坐标平面上所有的点组成的集合，如果由 \(A\) 到 \(A\) 的一一对应映射 \(f\)，使象集合的元素 \((y-1,x+2)\) 和原象集合的元素 \((x,y)\) 对应，那么，象点 \((3,-4)\) 的原象是（\quad）

\choices
    {\((-5,5)\)}
    {\((-6,4)\)}
    {\((2,-2)\)}
    {\((4,-6)\)}
\end{problem}

\begin{answer}
B
\end{answer}

\begin{solution}
映射把原象 \((x,y)\) 变为象 \((y-1,x+2)\). 若象点为 \((3,-4)\)，则
\(y-1=3\)，\(x+2=-4\)，从而 \(y=4\)，\(x=-6\). 所以原象点为 \((-6,4)\)，故选 B.
\end{solution}

\begin{problem}
不等式 \((x^2-4x-5)(x^2+8)<0\) 的解集是（\quad）

\choices
    {\(\{x\mid -1<x<5\}\)}
    {\(\{x\mid x<-1\text{\ 或\ }x>5\}\)}
    {\(\{x\mid 0<x<5\}\)}
    {\(\{x\mid -1<x<0\}\)}
\end{problem}

\begin{answer}
A
\end{answer}

\begin{solution}
因为 \(x^2+8>0\) 对一切实数 \(x\) 恒成立，所以原不等式等价于 \(x^2-4x-5<0\). 因式分解得 \((x-5)(x+1)<0\)，故 \(-1<x<5\)，解集为 \(\{x\mid-1<x<5\}\)，选 A.
\end{solution}

\begin{problem}
等差数列 \(\{a_n\}\) 的前三项 \(a_1=1\)，公差 \(d\neq0\)，如果 \(a_1\)，\(a_2\)，\(a_5\) 成等比数列，那么 \(d=\)（\quad）

\choices
    {\(3\)}
    {\(2\)}
    {\(-2\)}
    {\(2\)或\(-2\)}
\end{problem}

\begin{answer}
B
\end{answer}

\begin{solution}
由等差数列通项公式，\(a_2=1+d\)，\(a_5=1+4d\). 三项 \(a_1\)，\(a_2\)，\(a_5\) 成等比数列，所以 \(a_2^2=a_1a_5\)，即 \((1+d)^2=1+4d\). 化简得 \(d(d-2)=0\). 由于 \(d\neq0\)，故 \(d=2\)，选 B.
\end{solution}

\begin{problem}
函数 \(y=|x-1|\) 的图象是（\quad）

\choices
    {\choicebitmap[width=0.15\paperwidth]{img/1989/guangdong_liberal/q12_optA.png}}
    {\choicebitmap[width=0.15\paperwidth]{img/1989/guangdong_liberal/q12_optB.png}}
    {\choicebitmap[width=0.15\paperwidth]{img/1989/guangdong_liberal/q12_optC.png}}
    {\choicebitmap[width=0.15\paperwidth]{img/1989/guangdong_liberal/q12_optD.png}}
\end{problem}

\begin{answer}
B
\end{answer}

\begin{solution}
函数可写为
\[
y=|x-1|=\begin{cases}
1-x,&x<1,\\
x-1,&x\geqslant1.
\end{cases}
\]
它的图象是顶点为 \((1,0)\) 的开口向上的折线，并且在 \(x=0\) 时 \(y=1\)，与选项 B 的图象相符，故选 B.
\end{solution}

\begin{problem}
如果 \(a\)，\(b\)，\(c\) 都是正数，那么，在直角坐标系中，方程 \(ax+by+c=0\) 所表示的直线不经过（\quad）

\choices
    {第一象限}
    {第二象限}
    {第三象限}
    {第四象限}
\end{problem}

\begin{answer}
A
\end{answer}

\begin{solution}
把方程写成 \(y=-\frac abx-\frac cb\). 它与 \(x\) 轴、\(y\) 轴的截距分别为 \(-\frac ca<0\) 和 \(-\frac cb<0\). 直线的两个截距都在负半轴上，因此它在第二、第三、第四象限有点，但不可能同时有 \(x>0\) 与 \(y>0\) 的点，故不经过第一象限，选 A.
\end{solution}

\begin{problem}
如果椭圆 \(4x^2+y^2=k\) 上两点间的最大距离是 \(8\)，那么 \(k=\)（\quad）

\choices
    {\(32\)}
    {\(16\)}
    {\(8\)}
    {\(4\)}
\end{problem}

\begin{answer}
B
\end{answer}

\begin{solution}
椭圆方程化为 \(\frac{x^2}{k/4}+\frac{y^2}{k}=1\)，其两个半轴长为 \(\frac{\sqrt{k}}2\) 和 \(\sqrt{k}\). 两点间的最大距离是长轴长 \(2\sqrt{k}\)，所以 \(2\sqrt{k}=8\)，得 \(k=16\)，选 B.
\end{solution}

\begin{problem}
如果双曲线经过点 \((6,\sqrt{3})\)，且它的两条渐近线方程是 \(y=\pm \frac{1}{3}x\)，那么这双曲线的方程是（\quad）

\choices
    {\(\frac{x^2}{36}-\frac{y^2}{4}=1\)}
    {\(\frac{x^2}{81}-\frac{y^2}{9}=1\)}
    {\(\frac{x^2}{9}-y^2=1\)}
    {\(\frac{x^2}{18}-\frac{y^2}{3}=1\)}
\end{problem}

\begin{answer}
C
\end{answer}

\begin{solution}
渐近线为 \(y=\pm\frac13x\). 若双曲线为竖轴型 \(\frac{y^2}{a^2}-\frac{x^2}{b^2}=1\)，则由渐近线得 \(\frac ab=\frac13\)，代入点 \((6,\sqrt3)\) 有 \(\frac3{a^2}-\frac{36}{9a^2}=-\frac1{a^2}=1\)，矛盾. 因此双曲线只能为横轴型，设其方程为 \(\frac{x^2}{a^2}-\frac{y^2}{b^2}=1\). 由渐近线得 \(\frac ba=\frac13\)，即 \(b^2=\frac{a^2}{9}\). 代入点 \((6,\sqrt3)\)，得 \(\frac{36}{a^2}-\frac{3}{a^2/9}=1\)，即 \(\frac9{a^2}=1\)，所以 \(a^2=9\)，\(b^2=1\). 所求方程为 \(\frac{x^2}{9}-y^2=1\)，选 C.
\end{solution}

\begin{problem}
如果 \(l\) 和 \(n\) 是异面直线，那么与 \(l\)，\(n\) 都垂直的直线（\quad）

\choices
    {不一定存在}
    {总共只有一条}
    {总共可能有一条，也可能有两条}
    {必然有无穷多条}
\end{problem}

\begin{answer}
D
\end{answer}

\begin{solution}
设 \(l\)、\(n\) 的方向向量分别为 \(\bs{u}\)、\(\bs{v}\). 异面直线不平行，所以非零向量 \(\bs{w}=\bs{u}\times\bs{v}\) 同时垂直于 \(\bs{u}\) 和 \(\bs{v}\). 因此所有方向为 \(\bs{w}\) 的直线都与 \(l\)、\(n\) 垂直；平移这条方向的直线可得到无穷多条，故选 D.
\end{solution}

\begin{problem}
侧面都是直角三角形的正三棱锥，当底面边长为 \(a\) 时，该三棱锥的全面积是（\quad）

\choices
    {\(\frac{3+\sqrt{3}}{4}a^2\)}
    {\(\frac{3}{4}a^2\)}
    {\(\frac{3+3\sqrt{3}}{2}a^2\)}
    {\(\frac{6+\sqrt{3}}{4}a^2\)}
\end{problem}

\begin{answer}
A
\end{answer}

\begin{solution}
正三棱锥的每个侧面都是等腰三角形；又因其为直角三角形，底面边长 \(a\) 必为该等腰直角三角形的斜边. 所以每个侧面面积为 \(\frac12\left(\frac a{\sqrt2}\right)^2=\frac{a^2}{4}\)，三个侧面面积为 \(\frac{3a^2}{4}\). 底面是边长为 \(a\) 的正三角形，面积为 \(\frac{\sqrt3}{4}a^2\)，故全面积为 \(\frac{3+\sqrt3}{4}a^2\)，选 A.
\end{solution}

\begin{problem}
如果 \(A\) 是第四象限的角，且 \(\cos A=\frac{4}{5}\)，那么 \(\tan \frac{A}{2}=\)（\quad）

\choices
    {\(-3\)}
    {\(\frac{1}{3}\)或\(-\frac{1}{3}\)}
    {\(\frac{1}{3}\)}
    {\(-\frac{1}{3}\)}
\end{problem}

\begin{answer}
D
\end{answer}

\begin{solution}
因为 \(A\) 是第四象限角，\(\sin A<0\). 由 \(\cos A=\frac45\) 和 \(\sin^2A+\cos^2A=1\)，得 \(\sin A=-\frac35\). 半角公式给出 \(\tan\frac A2=\frac{\sin A}{1+\cos A}=\frac{-3/5}{1+4/5}=-\frac13\)，故选 D.
\end{solution}

\begin{problem}
函数 \(y=\sin\left(x+\frac{\pi}{5}\right)\cdot\sin\left(x-\frac{\pi}{5}\right)\) 的最小正周期是（\quad）

\choices
    {\(4\pi\)}
    {\(2\pi\)}
    {\(\pi\)}
    {\(\frac{1}{2}\pi\)}
\end{problem}

\begin{answer}
C
\end{answer}

\begin{solution}
利用积化和差公式，
\[
y=\frac12\left[\cos\frac{2\pi}{5}-\cos(2x)\right]
\]
其中常数项不影响周期，\(\cos(2x)\) 的最小正周期为 \(\pi\)，且其系数不为零，所以原函数的最小正周期为 \(\pi\)，选 C.
\end{solution}

\begin{problem}
函数 \(y=2\sin^2 x+2\cos x-3\) 的最大值是（\quad）

\choices
    {\(-1\)}
    {\(-\frac{1}{2}\)}
    {\(\frac{1}{2}\)}
    {\(-5\)}
\end{problem}

\begin{answer}
B
\end{answer}

\begin{solution}
令 \(t=\cos x\)，则 \(-1\leqslant t\leqslant1\)，且
\[
y=2(1-t^2)+2t-3=-2t^2+2t-1=-2\left(t-\frac12\right)^2-\frac12
\]
当 \(t=\frac12\) 时等号成立，故最大值为 \(-\frac12\)，选 B.
\end{solution}

\begin{problem}
如果函数 \(f(x)=x^2+2(a-1)x+2\) 在区间 \((-\infty,4]\) 上是减函数，那么实数 \(a\) 的取值范围是（\quad）

\choices
    {\(a\leqslant -3\)}
    {\(a\geqslant -3\)}
    {\(a\leqslant 5\)}
    {\(a\geqslant 3\)}
\end{problem}

\begin{answer}
A
\end{answer}

\begin{solution}
配方得 \(f(x)=(x+a-1)^2+2-(a-1)^2\)，抛物线开口向上，顶点横坐标为 \(1-a\). 函数在 \((-\infty,1-a]\) 上递减，因此要在整个 \((-\infty,4]\) 上递减，必须且只需 \(4\leqslant1-a\)，即 \(a\leqslant-3\)，选 A.
\end{solution}

\begin{problem}
如果 \(x\)，\(y\) 是实数，那么 \(\cos x=\cos y\) 是“\(x=y\)”的（\quad）

\choices
    {必要而且充分的条件}
    {必要但不充分的条件}
    {充分但不必要的条件}
    {既不充分也不必要的条件}
\end{problem}

\begin{answer}
B
\end{answer}

\begin{solution}
若 \(x=y\)，显然有 \(\cos x=\cos y\)，所以 \(\cos x=\cos y\) 是 \(x=y\) 的必要条件. 但反过来不成立，例如 \(x=0\)、\(y=2\pi\) 时余弦值相等而 \(x\neq y\). 因此它是必要但不充分的条件，选 B.
\end{solution}

\begin{problem}
满足条件 \(|z|=1\) 及 \(\left|z+\frac{1}{2}\right|=\left|z-\frac{3}{2}\right|\) 的复数 \(z\) 的集合是（\quad）

\choices
    {\(\left\{-\frac{1}{2}+\frac{\sqrt{3}}{2}\i,-\frac{1}{2}-\frac{\sqrt{3}}{2}\i\right\}\)}
    {\(\left\{\frac{1}{2}+\frac{1}{2}\i,\frac{1}{2}-\frac{1}{2}\i\right\}\)}
    {\(\left\{\frac{\sqrt{2}}{2}-\frac{\sqrt{2}}{2}\i,\frac{\sqrt{2}}{2}+\frac{\sqrt{2}}{2}\i\right\}\)}
    {\(\left\{\frac{1}{2}+\frac{\sqrt{3}}{2}\i,\frac{1}{2}-\frac{\sqrt{3}}{2}\i\right\}\)}
\end{problem}

\begin{answer}
D
\end{answer}

\begin{solution}
设 \(z=u+v\i\)，其中 \(u\)，\(v\in\R\). 由 \(|z|=1\)，得 \(u^2+v^2=1\). 由两个模相等，得 \((u+\frac12)^2+v^2=(u-\frac32)^2+v^2\)，化简为 \(u=\frac12\). 代回第一式得 \(v^2=\frac34\)，所以 \(v=\pm\frac{\sqrt3}{2}\). 因而 \(z=\frac12\pm\frac{\sqrt3}{2}\i\)，选 D.
\end{solution}

\begin{problem}
如果 \(x\)，\(y\in\left(0,\frac{\pi}{2}\right)\)，且 \(\tan x<\cot y\)，那么（\quad）

\choices
    {\(x+y>\frac{\pi}{2}\)}
    {\(x+y<\frac{\pi}{2}\)}
    {\(x>y\)}
    {\(x<y\)}
\end{problem}

\begin{answer}
B
\end{answer}

\begin{solution}
在 \(\left(0,\frac\pi2\right)\) 上，\(\cot y=\tan\left(\frac\pi2-y\right)\)，且正切函数严格递增. 所以 \(\tan x<\tan\left(\frac\pi2-y\right)\) 等价于 \(x<\frac\pi2-y\)，即 \(x+y<\frac\pi2\)，故选 B.
\end{solution}

\begin{problem}
过抛物线 \(y^2=4x\) 的焦点作直线交抛物线于 \(A(x_1,y_1)\)，\(B(x_2,y_2)\) 两点，如果 \(x_1+x_2=6\)，那么 \(|AB|=\)（\quad）

\choices
    {\(10\)}
    {\(8\)}
    {\(6\)}
    {\(4\)}
\end{problem}

\begin{answer}
B
\end{answer}

\begin{solution}
设抛物线上的点用参数表示为 \(P(t)=(t^2,2t)\)，记 \(A=P(t_1)\)、\(B=P(t_2)\). 焦点为 \(F(1,0)\). 弦 \(AB\) 过焦点的条件为 \(t_1t_2=-1\)：事实上，弦的方程为 \(y=\frac{2}{t_1+t_2}x+\frac{2t_1t_2}{t_1+t_2}\)，代入 \((1,0)\) 即得该条件；竖直弦不满足题设的 \(x_1+x_2=6\)，故参数式适用.

由 \(x_1+x_2=t_1^2+t_2^2=6\)，得 \((t_1+t_2)^2-2t_1t_2=6\)，从而 \((t_1+t_2)^2=4\). 又
\[
|AB|^2=(t_1^2-t_2^2)^2+4(t_1-t_2)^2=(t_1-t_2)^2\bigl((t_1+t_2)^2+4\bigr)
\]
其中 \((t_1-t_2)^2=(t_1+t_2)^2-4t_1t_2=8\)，所以 \(|AB|^2=8\cdot8=64\)，得 \(|AB|=8\)，选 B.
\end{solution}

\section{填空题}

\begin{problem}

\(\tan 315^\circ - \tan(-300^\circ) + \cot(-330^\circ)\) 的值是\fillinblank{}.
\end{problem}

\begin{answer}
\(-1\)
\end{answer}

\begin{solution}
利用周期性，\(\tan315^\circ=\tan(360^\circ-45^\circ)=-1\)，\(\tan(-300^\circ)=\tan60^\circ=\sqrt3\)，\(\cot(-330^\circ)=\cot30^\circ=\sqrt3\). 因而原式 \(=-1-\sqrt3+\sqrt3=-1\).
\end{solution}

\begin{problem}

函数 \(y = \frac{1}{3x + 1}\) 的反函数是\fillinblank{}.
\end{problem}

\begin{answer}
\(y = \frac{1 - x}{3x}\)
\end{answer}

\begin{solution}
设原函数为 \(y=\frac1{3x+1}\)，则 \(y\neq0\). 交换 \(x\)，\(y\) 得 \(x=\frac1{3y+1}\)，所以 \(x(3y+1)=1\)，即 \(3xy=1-x\). 因而反函数为 \(y=\frac{1-x}{3x}\)，其定义域为 \(x\neq0\).
\end{solution}

\begin{problem}

过圆 \(x^2 + y^2 = 12\) 上的点 \(M(3,\sqrt{3})\) 作圆的切线，这切线的方程是\fillinblank{}.
\end{problem}

\begin{answer}
\(3x + \sqrt{3}y = 12\)
\end{answer}

\begin{solution}
圆心为原点，半径平方为 \(12\). 半径 \(OM\) 的方向向量为 \((3,\sqrt3)\)，切线过 \(M(3,\sqrt3)\) 且与 \(OM\) 垂直，因此切线方程为 \(3x+\sqrt3y=3^2+(\sqrt3)^2=12\).
\end{solution}

\begin{problem}

\(\displaystyle\lim_{n \to \infty}\left(\frac{1}{n^2 + 1} + \frac{2}{n^2 + 1} + \cdots + \frac{n}{n^2 + 1}\right)\) 的值是\fillinblank{}.
\end{problem}

\begin{answer}
\(\frac{1}{2}\)
\end{answer}

\begin{solution}
把各项合并，得
\[
\frac{1}{n^2+1}+\frac{2}{n^2+1}+\cdots+\frac{n}{n^2+1}=\frac{1+2+\cdots+n}{n^2+1}=\frac{n(n+1)}{2(n^2+1)}
\]
分子分母同除以 \(n^2\)，当 \(n\to\infty\) 时极限为 \(\frac12\).
\end{solution}

\begin{problem}

由 \(1\)，\(2\)，\(3\)，\(4\) 四个数字组成的没有重复数字的四位数中，比 \(1234\) 大的共有\fillinblank{}个（用具体数字作答）.
\end{problem}

\begin{answer}
23
\end{answer}

\begin{solution}
按千位数字分类. 千位为 \(2\)、\(3\)、\(4\) 时，后三位分别有 \(3!=6\) 个，共 \(18\) 个. 千位为 \(1\) 时，百位为 \(2\) 时只有 \(1243\) 大于 \(1234\)，有 \(1\) 个；百位为 \(3\) 或 \(4\) 时，后两位各有 \(2!=2\) 个，共 \(5\) 个. 因而总数为 \(18+5=23\) 个.
\end{solution}

\section{解答题}

\begin{problem}

设 \(f(x) = a^{x - 1}\)，其中常数 \(a > 0\).

\begin{enumerate}
\item 如果 \(\{x_n\}\) 是等差数列，求证：\(\{f(x_n)\}\) 是等比数列；
\item 如果 \(a = 4\)，解方程 \(4f(x) + f(-x) - 1 = 0\).
\end{enumerate}
\end{problem}

\begin{solution}
\begin{enumerate}
\item 设等差数列 \(\{x_n\}\) 的公差为 \(d\)，则 \(x_n-x_{n-1}=d\). 因为 \(a>0\)，所以每一项 \(f(x_n)=a^{x_n-1}>0\). 当 \(n\geqslant2\) 时，
\[
\frac{f(x_n)}{f(x_{n-1})}=\frac{a^{x_n-1}}{a^{x_{n-1}-1}}=a^d
\]
右端是与 \(n\) 无关且为正的常数，所以 \(\{f(x_n)\}\) 是等比数列，公比为 \(a^d\).

\item 当 \(a=4\) 时，\(4f(x)=4\cdot4^{x-1}=4^x\)，且 \(f(-x)=4^{-x-1}\)，所以方程化为 \(4^x+4^{-x-1}-1=0\). 令 \(y=4^x>0\)，则
\[
y+\frac{1}{4y}-1=0
\]
乘以正数 \(4y\)，得 \(4y^2-4y+1=0\)，即 \((2y-1)^2=0\)，所以 \(y=\frac12\). 因而 \(4^x=\frac12=4^{-1/2}\)，解得 \(x=-\frac12\).
\end{enumerate}
\end{solution}

\begin{problem}

如图，在平面 \(\beta\) 内有 \(\triangle ABC\)，在平面 \(\beta\) 外有点 \(S\)，斜线 \(SA \perp AC\)，\(SB \perp BC\)，且斜线段 \(SA\) 和 \(SB\) 在 \(\beta\) 内的射影是等长的线段.

\begin{enumerate}
\item 求证：\(AC = BC\)；
\item 又设点 \(S\) 与平面 \(\beta\) 的距离为 \(3\mathrm{~cm}\)，且 \(AC \perp BC\)，\(AB = 4\mathrm{~cm}\)，求线段 \(SC\) 的长.
\end{enumerate}

\bitmapfigure[max width=0.42\linewidth]{img/1989/guangdong_liberal/q2_fig1.png}
\end{problem}

\begin{solution}
\begin{enumerate}
\item 作 \(SD\perp\beta\) 于 \(D\)，连接 \(AD\)、\(BD\). 斜线段 \(SA\)、\(SB\) 在 \(\beta\) 内的射影分别为 \(AD\)、\(BD\)，所以 \(AD=BD\). 在直角三角形 \(SAD\)、\(SBD\) 中，\(SD\) 是公共直角边，故 \(SA^2=SD^2+AD^2=SD^2+BD^2=SB^2\)，从而 \(SA=SB\). 又 \(SA\perp AC\)、\(SB\perp BC\)，所以 \(\triangle SAC\)、\(\triangle SBC\) 都是直角三角形，且有公共斜边 \(SC\) 和另一条直角边 \(SA=SB\)，由直角三角形全等得 \(AC=BC\).
\item 因为 \(AC\perp SA\)、\(AC\subset\beta\)，由三垂线定理的逆定理得 \(AC\perp AD\). 又因为 \(SB\) 在平面 \(\beta\) 内的射影是 \(BD\)，且 \(SB\perp BC\)、\(BC\subset\beta\)，由同一定理的逆定理得 \(BC\perp BD\)，即 \(\angle DAC=\angle DBC=90^\circ\). 再由 \(AC\perp BC\) 和 \(AC=BC\)，在平面 \(\beta\) 内取 \(C\) 为原点、\(CA\)、\(CB\) 为两条互相垂直的坐标轴，可知 \(D\) 的两个坐标分别等于 \(CA\)、\(CB\)，故四边形 \(ACBD\) 为正方形. 因而 \(CD=AB=4\mathrm{~cm}\). 又 \(SD\perp\beta\)，所以 \(SD\perp CD\)，在直角三角形 \(SCD\) 中，\(SC=\sqrt{SD^2+CD^2}=\sqrt{3^2+4^2}=5\mathrm{~cm}\).
\end{enumerate}

\solutionfigure{\bitmapfigure[max width=0.42\linewidth]{img/1989/guangdong_liberal/q2_solution_fig1.png}}
\end{solution}

\begin{problem}

已知复数 \(z = 1 - 2\i\)，其中 \(\i\) 是虚数单位.

\begin{enumerate}
\item 求 \(\frac{2}{z + \i} - \overline{(z + 1)}\) 的值，并把该值表示成三角形式；
\item 求适合不等式 \(\log_{0.5}\left|\frac{az - \i}{\sqrt{a + 1}}\right| \leqslant \frac{1}{2}\) 的实数 \(a\) 的取值范围.
\end{enumerate}
\end{problem}

\begin{solution}
\begin{enumerate}
\item 因为 \(z = 1 - 2\i\)，所以
\[
\frac{2}{z + \i} - \overline{\left(z + 1\right)}
= \frac{2}{1 - \i} - \overline{\left(2 - 2\i\right)}
= 1 + \i - \left(2 + 2\i\right)
= -1 - \i
\]
将其表示成三角形式，得 \(-1 - \i = \sqrt{2}\left(\cos \frac{5\pi}{4} + \i \sin \frac{5\pi}{4}\right)\).

\item 因为 \(z=1-2\i\)，所以 \(az-\i=a-(2a+1)\i\)，从而 \(|az-\i|=\sqrt{a^2+(2a+1)^2}=\sqrt{5a^2+4a+1}\). 对数的真数必须为正，且分母中的根式要求 \(a+1>0\). 在此条件下真数为正，因为 \(5a^2+4a+1=a^2+(2a+1)^2\) 不可能为零. 由于底数 \(0.5\) 在 \(0\) 与 \(1\) 之间，原不等式等价于
\[
\frac{\sqrt{5a^2+4a+1}}{\sqrt{a+1}}\geqslant\left(\frac12\right)^{1/2}
\]
两边均非负，可以平方并利用 \(a+1>0\)，得
\[
\frac{5a^2+4a+1}{a+1}\geqslant\frac12
\]
即 \(10a^2+7a+1\geqslant0\). 因式分解为 \((5a+1)(2a+1)\geqslant0\)，所以 \(a\leqslant-\frac12\) 或 \(a\geqslant-\frac15\). 再与定义域 \(a>-1\) 取交集，得到 \(a\in\left(-1,-\frac12\right]\cup\left[-\frac15,+\infty\right)\).
\end{enumerate}
\end{solution}

\begin{problem}

设圆 \(O\) 的方程为 \(x^2 - 2x\cos \theta + y^2 + 2(1 + 2\sin \theta)y = 0\)，其中 \(0\leqslant\theta<2\pi\).

\begin{enumerate}
\item 当 \(\theta\) 变化时，求圆 \(O\) 的圆心轨迹的普通方程，并画出轨迹的草图；
\item 如果又有圆 \(S\)：\(x^2 + y^2 + 2y = 0\)，求 \(\theta\) 的值，使得圆 \(O\) 与圆 \(S\) 的公共弦的长度取最大值（请把满足条件的 \(\theta\) 的值都求出来）.
\end{enumerate}
\end{problem}

\begin{solution}
\begin{enumerate}
\item 配方可知圆 \(O\) 的圆心为 \(C\left(\cos\theta,-1-2\sin\theta\right)\). 令圆心坐标为 \((x,y)\)，则 \(x=\cos\theta\)，\(y+1=-2\sin\theta\)，消去 \(\theta\) 得
\[
x^2+\frac{(y+1)^2}{4}=1
\]
即
\[
4x^2+(y+1)^2=4
\]
因此圆心轨迹是中心为 \((0,-1)\)，横、纵半轴长分别为 \(1\) 和 \(2\) 的椭圆. 草图可由四个顶点 \((1,-1)\)、\((-1,-1)\)、\((0,1)\)、\((0,-3)\) 确定，并以 \((0,-1)\) 为中心作长轴竖直的椭圆.

\item 两圆方程相减，得公共弦所在直线 \(x\cos\theta-2y\sin\theta=0\). 圆 \(S\) 的圆心为 \(S_0(0,-1)\)，半径为 \(1\)，圆心到该直线的距离为
\[
p=\frac{2|\sin\theta|}{\sqrt{\cos^2\theta+4\sin^2\theta}}=\frac{2|\sin\theta|}{\sqrt{1+3\sin^2\theta}}
\]
故公共弦是圆 \(S\) 中心距圆心距离为 \(p\) 的弦，其长度为
\[
L=2\sqrt{1-p^2}=\frac{2|\cos\theta|}{\sqrt{1+3\sin^2\theta}}
\]
因为 \(0\leqslant\sin^2\theta\leqslant1\)，有 \(L^2=4\frac{1-\sin^2\theta}{1+3\sin^2\theta}\leqslant4\)，且等号当且仅当 \(\sin\theta=0\). 所以公共弦长度最大值为 \(2\)，由 \(0\leqslant\theta<2\pi\)，得 \(\theta=0\) 或 \(\pi\).
\end{enumerate}
\end{solution}
